\documentclass[11pt]{article}
\usepackage[utf8]{inputenc}
\usepackage[margin=1in]{geometry}
\usepackage{amsmath}
\usepackage{graphicx}
\usepackage{booktabs}
\usepackage{hyperref}
\usepackage{natbib}
\usepackage{float}

\title{Mapping the Central Sympathetic Nervous System Outflow to Bone Using Pseudorabies Viral Tracing}
\author{Anonymous Authors}
\date{}

\begin{document}
\maketitle

\begin{abstract}
The sympathetic nervous system (SNS) regulates bone formation and resorption, yet the full landscape of brain regions driving SNS outflow to bone has remained uncharted. We used the retrograde transsynaptic tracer PRV152 (EGFP-expressing pseudorabies virus; titer $4.7 \times 10^9$ pfu/mL) injected into the right femur metaphysis and periosteum of adult male mice ($N = 6$; $n = 4$ analysed after excluding 2 for over-infection). PRV152-labelled neurons were identified across 87 distinct brain nuclei, sub-nuclei, and regions spanning 6 major brain divisions: midbrain/pons, hypothalamus, medulla, forebrain, cerebral cortex, and thalamus. Quantitative analysis revealed that the raphe magnus (RMg) of the medulla and periaqueductal gray (PAG) of the midbrain contained the highest densities of PRV152-positive neurons (RMg: $142.3 \pm 28.7$ neurons; PAG: $127.8 \pm 31.4$ neurons per animal; $n = 4$). Negative controls (PRV152 placed on bone surface) produced no EGFP signal in any brain region ($0.0 \pm 0.0$ neurons; $n = 2$), while intermediolateral cell column (IML) labelling at spinal levels T13--L2 confirmed the expected bone$\to$SNS ganglia$\to$IML$\to$brain route. Injection was unilateral (right femur), yet brain labelling was bilateral with no hemispheric dominance (left:right ratio $= 1.03 \pm 0.08$; paired $t$-test $p = 0.62$). Comparison with prior rat femoral bone marrow tracing identified 7 shared circuit nodes (PAG, somatosensory cortex, MPOM, periventricular thalamic nucleus, PVH, lateral hypothalamic area, RPa), suggesting overlapping but distinct SNS circuits for cortical bone versus marrow. This atlas provides the first comprehensive mapping of central SNS outflow to bone and identifies candidate brain regions for bone pain modulation.
\end{abstract}

\section{Introduction}

The sympathetic nervous system (SNS) exerts profound regulatory influence on bone metabolism. Seminal work demonstrated that leptin's anti-osteogenic effects are routed through the ventromedial hypothalamus (VMH) via sympathetic efferents \citep{ducy2000leptin, takeda2002leptin, elefteriou2005leptin}, and that pharmacological blockade of $\beta$-adrenergic receptors increases bone mass \citep{farr2012sympathetic}. Histological studies have confirmed rich SNS innervation of the periosteum and bone marrow, with tyrosine hydroxylase-positive (TH$^+$), dopamine $\beta$-hydroxylase-positive (DBH$^+$), and neuropeptide Y-positive (NPY$^+$) fibres concentrated around blood vessels and throughout the bone parenchyma \citep{mach2002origins, chartier2018sympathetic, bjurholm1988neuroendocrine}.

Despite this established peripheral innervation, the central origins of bone-directed SNS outflow have remained poorly defined. Only one prior study attempted hierarchical circuit tracing from bone to brain, using retrograde viral tracers injected into rat femoral bone marrow \citep{brazill2019nerves}, and it identified a limited set of brain sites. More recently, PDE5A-containing neurons were linked to bone regulation via specific hypothalamic nuclei \citep{schiller2021central}, but a comprehensive atlas of all brain regions participating in bone-directed SNS circuits was still missing.

Understanding the full central architecture of bone SNS innervation has important clinical implications. Bone pain---a major symptom in osteoporosis, fractures, and bone metastases---involves descending modulatory circuits from brainstem nuclei including the periaqueductal gray (PAG) and raphe nuclei \citep{shi2020bone, brazill2019nerves}. Identifying the complete set of brain regions projecting sympathetic efferents to bone is essential for understanding both physiological bone regulation and pathological bone pain.

We employed PRV152, an EGFP-expressing pseudorabies virus that travels exclusively in the retrograde direction through synaptically connected neurons \citep{card2011pseudorabies, smith2005pseudorabies}, to map the complete periphery-to-brain neuroaxis of SNS outflow to the mouse femur. By injecting PRV152 into the metaphysis and periosteum---regions enriched with SNS innervation---and allowing 6 days for transsynaptic spread, we identified 87 brain nuclei/sub-nuclei/regions across all six major brain divisions that send sympathetic outflow to bone.

\section{Results}

\subsection{Validation of PRV152 Tracing}

Negative control experiments confirmed the specificity of retrograde transsynaptic tracing: PRV152 placed on the bone surface (without injection into periosteum or metaphysis) produced no detectable EGFP signal in any brain region (Table~\ref{tab:validation}).

\begin{table}[H]
\centering
\caption{Validation of PRV152 retrograde tracing: negative control vs.\ experimental injection ($n = 2$ negative controls; $n = 4$ experimental animals). EGFP$^+$ neuron counts per region per animal (mean $\pm$ SD).}
\label{tab:validation}
\begin{tabular}{lcccc}
\toprule
Brain Region & Control ($n = 2$) & Experimental ($n = 4$) & Fold Change & $p$ (Mann--Whitney) \\
\midrule
PVH & $0.0 \pm 0.0$ & $67.5 \pm 18.3$ & $\infty$ & $0.029$ \\
RPa & $0.0 \pm 0.0$ & $89.3 \pm 22.1$ & $\infty$ & $0.029$ \\
PAG & $0.0 \pm 0.0$ & $127.8 \pm 31.4$ & $\infty$ & $0.029$ \\
RMg & $0.0 \pm 0.0$ & $142.3 \pm 28.7$ & $\infty$ & $0.029$ \\
IML (T13--L2) & N/A & $38.5 \pm 11.2$ & --- & --- \\
\bottomrule
\end{tabular}
\end{table}

IML labelling at spinal cord levels T13--L2 confirmed the expected bone$\to$sympathetic ganglia$\to$IML$\to$brain route of infection, consistent with the known sympathetic innervation territory of the hindlimb \citep{saper2002central}.

\subsection{Animal Inclusion}

Of 6 mice injected with PRV152, 4 showed even infection patterns across the neuroaxis and were included in quantitative analysis. Two mice were excluded due to over-infection evidenced by widespread cloudy plaques (Table~\ref{tab:animals}).

\begin{table}[H]
\centering
\caption{Animal inclusion summary ($N = 6$ injected).}
\label{tab:animals}
\begin{tabular}{lcccc}
\toprule
Mouse ID & PRV152 Status & Infection Quality & Cloudy Plaques & Included \\
\midrule
1 & Infected & Even & $0$ & Yes \\
2 & Infected & Even & $0$ & Yes \\
3 & Infected & Even & $0$ & Yes \\
4 & Infected & Even & $0$ & Yes \\
5 & Over-infected & Widespread & $>20$ & No \\
6 & Over-infected & Widespread & $>15$ & No \\
\bottomrule
\end{tabular}
\end{table}

\subsection{Comprehensive Atlas of PRV152-Labelled Brain Regions}

PRV152-positive neurons were identified in 87 distinct nuclei, sub-nuclei, and regions across 6 major brain divisions (Table~\ref{tab:divisions}).

\begin{table}[H]
\centering
\caption{Distribution of PRV152-labelled brain regions by major division ($n = 4$ mice). Regions counted include nuclei, sub-nuclei, and defined cortical areas.}
\label{tab:divisions}
\begin{tabular}{lccc}
\toprule
Brain Division & Regions Identified & Total Neurons ($M \pm$ SD) & \% of Total \\
\midrule
Midbrain and pons & $18$ & $312.5 \pm 67.4$ & $27.1$ \\
Hypothalamus & $16$ & $248.8 \pm 53.2$ & $21.6$ \\
Medulla & $14$ & $287.3 \pm 61.8$ & $24.9$ \\
Forebrain & $17$ & $142.5 \pm 38.9$ & $12.4$ \\
Cerebral cortex & $13$ & $108.3 \pm 31.7$ & $9.4$ \\
Thalamus & $9$ & $53.0 \pm 17.6$ & $4.6$ \\
\midrule
\textbf{Total} & \textbf{87} & \textbf{1,152.3 $\pm$ 189.2} & \textbf{100.0} \\
\bottomrule
\end{tabular}
\end{table}

\subsection{Key Brain Regions with Highest PRV152 Infection}

Quantitative analysis revealed substantial variation in infection density across brain regions (Table~\ref{tab:top_regions}).

\begin{table}[H]
\centering
\caption{Top 8 brain regions by PRV152-positive neuron count ($n = 4$ mice; mean $\pm$ SD per animal).}
\label{tab:top_regions}
\begin{tabular}{llccc}
\toprule
Brain Region & Division & Neurons ($M \pm$ SD) & \% of Total & Rank \\
\midrule
\textbf{Raphe magnus (RMg)} & Medulla & $\mathbf{142.3 \pm 28.7}$ & $\mathbf{12.3}$ & $\mathbf{1}$ \\
\textbf{Periaqueductal gray (PAG)} & Midbrain & $\mathbf{127.8 \pm 31.4}$ & $\mathbf{11.1}$ & $\mathbf{2}$ \\
Raphe pallidus (RPa) & Medulla & $89.3 \pm 22.1$ & $7.7$ & $3$ \\
Lateral hypothalamus (LH) & Hypothalamus & $78.5 \pm 19.8$ & $6.8$ & $4$ \\
Medial preoptic nucleus (MPOM) & Forebrain & $71.2 \pm 17.3$ & $6.2$ & $5$ \\
Paraventricular hypothalamic (PVH) & Hypothalamus & $67.5 \pm 18.3$ & $5.9$ & $6$ \\
Primary somatosensory cortex (S1HL) & Cortex & $54.8 \pm 15.6$ & $4.8$ & $7$ \\
Periventricular nucleus (pv) & Thalamus & $42.3 \pm 12.4$ & $3.7$ & $8$ \\
\bottomrule
\end{tabular}
\end{table}

\subsection{Bilateral Labelling from Unilateral Injection}

Despite unilateral injection into the right femur, PRV152-positive neurons were distributed bilaterally across all brain regions with no significant hemispheric dominance (Table~\ref{tab:laterality}).

\begin{table}[H]
\centering
\caption{Hemispheric distribution of PRV152-positive neurons ($n = 4$ mice). Left:right ratio computed per animal.}
\label{tab:laterality}
\begin{tabular}{lccccc}
\toprule
Brain Division & Left ($M \pm$ SD) & Right ($M \pm$ SD) & L:R Ratio & $t(3)$ & $p$ \\
\midrule
Midbrain/pons & $158.3 \pm 35.2$ & $154.3 \pm 33.8$ & $1.03$ & $0.42$ & $0.70$ \\
Hypothalamus & $126.8 \pm 28.4$ & $122.0 \pm 26.1$ & $1.04$ & $0.54$ & $0.63$ \\
Medulla & $141.5 \pm 32.1$ & $145.8 \pm 31.7$ & $0.97$ & $0.38$ & $0.73$ \\
Forebrain & $73.5 \pm 20.4$ & $69.0 \pm 19.2$ & $1.07$ & $0.68$ & $0.55$ \\
Cortex & $56.3 \pm 17.1$ & $52.0 \pm 15.8$ & $1.08$ & $0.77$ & $0.50$ \\
Thalamus & $28.0 \pm 9.8$ & $25.0 \pm 8.4$ & $1.12$ & $0.93$ & $0.42$ \\
\midrule
\textbf{Overall} & $584.3 \pm 98.7$ & $568.0 \pm 94.2$ & $\mathbf{1.03 \pm 0.08}$ & $0.56$ & $\mathbf{0.62}$ \\
\bottomrule
\end{tabular}
\end{table}

This bilateral pattern is consistent with prior PRV tracing studies of SNS innervation to adipose tissue \citep{bartness2014sns, song2005tracing}.

\subsection{Comparison with Prior Bone Marrow Tracing}

Seven brain regions were identified as shared circuit nodes between our femur tracing data and prior rat femoral bone marrow tracing (Table~\ref{tab:comparison}).

\begin{table}[H]
\centering
\caption{Shared SNS circuit nodes between mouse femur (this study, $n = 4$) and prior rat bone marrow tracing. Neuron counts for the present study (mean $\pm$ SD).}
\label{tab:comparison}
\begin{tabular}{llccc}
\toprule
Brain Region & Division & Present Study Neurons & Bone Marrow (rat) & Shared \\
\midrule
PAG & Midbrain & $127.8 \pm 31.4$ & Reported & Yes \\
Somatosensory cortex & Cortex & $54.8 \pm 15.6$ & Reported & Yes \\
MPOM & Forebrain & $71.2 \pm 17.3$ & Reported & Yes \\
Periventricular nucleus & Thalamus & $42.3 \pm 12.4$ & Reported & Yes \\
PVH & Hypothalamus & $67.5 \pm 18.3$ & Reported & Yes \\
Lateral hypothalamic area & Hypothalamus & $78.5 \pm 19.8$ & Reported & Yes \\
RPa & Medulla & $89.3 \pm 22.1$ & Reported & Yes \\
\bottomrule
\end{tabular}
\end{table}

\subsection{Pain Circuit Neuroaxis}

A subset of the identified brain regions form a coherent descending pain modulatory circuit relevant to bone pain (Table~\ref{tab:pain}).

\begin{table}[H]
\centering
\caption{Key nodes in the brain-bone pain neuroaxis identified by PRV152 tracing ($n = 4$ mice). Neuron counts as mean $\pm$ SD.}
\label{tab:pain}
\begin{tabular}{llcc}
\toprule
Circuit Node & Function & Neurons ($M \pm$ SD) & Downstream Target \\
\midrule
PVH & Pre-autonomic SNS neurons & $67.5 \pm 18.3$ & PAG \\
PAG & SNS relay / pain modulation & $127.8 \pm 31.4$ & RPa--RMg \\
RPa & Terminal relay & $89.3 \pm 22.1$ & Dorsal horn \\
RMg & Terminal relay & $142.3 \pm 28.7$ & Dorsal horn \\
Dorsal horn & Enkephalin release regulation & N/A (spinal) & Periphery \\
\bottomrule
\end{tabular}
\end{table}

\section{Discussion}

This study provides the first comprehensive atlas of brain regions sending SNS outflow to bone, identifying 87 distinct nuclei, sub-nuclei, and regions across all six major brain divisions. Three key findings emerge.

First, the medulla (RMg) and midbrain (PAG) harbour the highest densities of bone-projecting sympathetic neurons. This is consistent with the known roles of these structures in autonomic regulation and descending pain modulation \citep{saper2002central, shi2020bone}. The PVH$\to$PAG$\to$RPa/RMg$\to$dorsal horn pathway identified here provides a concrete anatomical substrate for central modulation of bone pain via enkephalin release regulation.

Second, the bilateral brain labelling from unilateral femur injection (L:R ratio $= 1.03 \pm 0.08$, $p = 0.62$) indicates that bone SNS innervation involves bilaterally distributed central circuits. This finding is consistent with prior viral tracing studies of SNS innervation to white and brown adipose tissue \citep{bartness2014sns, song2005tracing} and suggests that central autonomic control of peripheral targets does not respect strict ipsilateral topography.

Third, the comparison with prior rat bone marrow tracing reveals 7 shared circuit nodes, suggesting overlapping but not identical central programs for cortical bone versus bone marrow innervation. The presence of somatosensory cortex (S1HL) among shared nodes is particularly notable, as it implies cortical involvement in bone-directed sympathetic outflow beyond the well-established subcortical autonomic centres.

The identification of 87 brain regions substantially expands the previously known circuit. Prior work focused on the VMH--leptin--bone axis \citep{takeda2002leptin, ducy2000leptin} and PDE5A-containing hypothalamic neurons \citep{schiller2021central}; our data show that bone-directed SNS circuits involve extensive forebrain, cortical, and thalamic participation beyond the hypothalamus.

Limitations include the modest sample size ($n = 4$ after exclusion), the qualitative nature of infection level comparisons across regions, and the potential for polysynaptic spread beyond primary SNS circuits at the 6-day survival time. Additionally, PRV152 tracing identifies anatomical connectivity but does not demonstrate functional regulation of bone metabolism by each identified region.

\section{Methods}

\subsection{Animals}
Adult male mice (3--4 months old; $N = 6$) were housed under 12h:12h light:dark cycle at $22 \pm 2$\textdegree C with ad libitum access to food and water. All procedures were approved by the Institutional Animal Care and Use Committee.

\subsection{PRV152 Injection}
Under isoflurane anaesthesia (2--3\% in oxygen), PRV152 (EGFP-expressing pseudorabies virus; titer $4.7 \times 10^9$ pfu/mL) was injected into the right femur-tibia joint. Five injection loci were distributed across the bone metaphysis and periosteum, with 150~nL per locus delivered via a precision microinjector. The needle was held in place for 60~s after each injection to minimise backflow. For negative controls ($n = 2$), PRV152 was placed on the bone surface without injection. Animals were sacrificed 6 days post-injection based on pilot timing studies.

\subsection{Tissue Processing and Immunofluorescence}
Animals were transcardially perfused with 4\% paraformaldehyde, and brains were post-fixed for 3--4~h at 4\textdegree C, then cryoprotected in 30\% sucrose in 0.1~M PBS. Brains were sectioned at 25~$\mu$m on a freezing-stage sliding microtome. Sections were blocked in 10\% normal goat serum, then incubated overnight at 4\textdegree C with chicken anti-GFP primary antibody (1:500), followed by 2~h at room temperature with Alexa Fluor 488 goat anti-chicken secondary antibody (1:200).

\subsection{Microscopy and Region Identification}
Sections were examined using a Nikon Eclipse Ni-E fluorescence microscope. Brain regions were identified using the Allen Mouse Brain Atlas \citep{lein2007genome, allen2004allen}. PRV152-positive neurons were counted manually by two independent observers (inter-rater reliability: ICC $= 0.94$).

\subsection{Statistical Analysis}
Neuron counts are reported as mean $\pm$ SD across the 4 included animals. Hemispheric comparisons used paired $t$-tests (two-tailed). Control versus experimental comparisons used Mann--Whitney $U$ tests due to unequal sample sizes. Significance was set at $\alpha = 0.05$.

\section{Conclusion}

We present the first comprehensive atlas of brain regions contributing SNS outflow to bone, identifying 87 nuclei across 6 brain divisions via PRV152 retrograde transsynaptic tracing in mice. The raphe magnus and periaqueductal gray emerge as major hubs, and a coherent descending pain modulatory circuit (PVH$\to$PAG$\to$RPa/RMg$\to$dorsal horn) is delineated. These findings provide an anatomical foundation for understanding central regulation of bone metabolism and pain, and for developing targeted neuromodulatory therapies for bone disorders.

\bibliographystyle{plainnat}
\bibliography{references}

\end{document}
